
\subsubsection{2.1 Communication Accommodation Theory}

Communication Accommodation Theory (CAT) (Giles, 1973; Giles and Ogay, 2007) posits that interlocutors converge or diverge linguistically as a function of social and relational factors. Danescu-Niculescu-Mizil et al. (2012) operationalized linguistic coordination using function-word category frequencies in Wikipedia administrator discussions and U.S. Supreme Court oral arguments, finding that lower-power interlocutors accommodate more to higher-power ones. Their coordination metric ($C_m$) operates at the lexical level over closed-class word categories. The present work extends accommodation measurement to continuous semantic embedding space (SBERT cosine similarity) and applies it to human-AI conversations stratified by RLHF quality tier.

\subsubsection{2.2 RLHF Alignment and the HH-RLHF Dataset}

RLHF (Christiano et al., 2017; Ouyang et al., 2022) uses human preference annotations to train reward models that guide language model optimization. Bai et al. (2022) introduced the HH-RLHF dataset, comprising three helpfulness tiers representing successive stages of alignment: supervised fine-tuning (helpful-base), rejection sampling with a reward model (helpful-rejection-sampled), and online PPO optimization (helpful-online). Existing work uses the tier structure to evaluate AI output quality. The present study uses the tier structure as an independent variable and measures human conversational behavior as the dependent variable.

\subsubsection{2.3 Human-AI Style Adaptation}

Chang and Wang (2025) demonstrated word-level bidirectional style adaptation in human-AI conversations across cultural contexts. Shen et al. (2025) proposed a bidirectional alignment framework motivating the study of mutual adaptation in human-AI interaction. Neither study examines SBERT-based semantic accommodation across RLHF quality tiers.

\subsubsection{2.4 Sentence Embeddings and Semantic Similarity}

Reimers and Gurevych (2019) introduced Sentence-BERT, producing semantically meaningful sentence embeddings via siamese BERT architectures with mean pooling. SBERT embeddings capture both content and style in a unified continuous space. This entanglement is a known limitation: cosine similarity between full-utterance embeddings does not cleanly separate topical from stylistic alignment. The three-level partner-specificity control hierarchy used in this work partially addresses this confound.
